\section{One-electron operator와 GTH pseudopotential derivative}
\label{sec:gth-one-electron}

\subsection{Periodic one-electron Hamiltonian}

GTH pseudopotential을 사용하는 one-electron Hamiltonian은
\begin{equation}
  \hat h
  =-\frac12\nabla^2
   +\sum_A\hat V_{\mathrm{loc}}^A
   +\sum_A\hat V_{\mathrm{nl}}^A
  \label{eq:gth-hamiltonian}
\end{equation}
이다. All-electron electron--nuclear $-Z_A/r$ 대신 smooth local potential과 finite-rank
nonlocal projector를 사용한다. Kinetic operator 자체는 핵 좌표에 명시적으로 의존하지
않지만 atom-centered AO basis가 움직이므로 kinetic AO matrix derivative는 0이 아니다.

\subsection{Local GTH potential}

대표적인 local form은
\begin{equation}
  V_{\mathrm{loc}}^A(r)
  =-\frac{Z_A^{\mathrm{ion}}}{r}
    \operatorname{erf}\!\left(\frac{r}{\sqrt2 r_{\mathrm{loc},A}}\right)
   +e^{-r^2/(2r_{\mathrm{loc},A}^2)}
    \sum_{n=1}^{4}C_{nA}
    \left(\frac{r}{r_{\mathrm{loc},A}}\right)^{2n-2}
  \label{eq:gth-local-real}
\end{equation}
이다. Periodic FFT representation에서는 원자 위치 dependence가 structure factor
$e^{-i\vG\cdot\mathbf R_A}$에 들어가므로
\begin{equation}
  \frac{\partial}{\partial R_{Ax}}
  \left[e^{-i\vG\cdot\mathbf R_A}V_{\mathrm{loc}}^A(\vG)\right]
  =-iG_xe^{-i\vG\cdot\mathbf R_A}V_{\mathrm{loc}}^A(\vG).
  \label{eq:gth-local-structure-derivative}
\end{equation}
$\vG=0$에서는 $G_x=0$이지만 constant term, charge convention과 Ewald compensation을
central hcore와 같은 방식으로 처리해야 한다.

AO matrix는
\begin{equation}
  V_{\mu\nu,\mathrm{loc}}^{\vk}
  =\langle\phi_{\mu\vk}|\hat V_{\mathrm{loc}}|\phi_{\nu\vk}\rangle
  \label{eq:gth-local-ao}
\end{equation}
이며 derivative에는 potential center $A$, bra AO center, ket AO center의 이동이 모두
포함된다. Translation identity를 사용하면 potential-center derivative를 AO derivative와
reciprocal structure-factor derivative의 조합으로 계산할 수 있다.

\subsection{Nonlocal projectors}

Nonlocal GTH operator는 angular momentum channel별 separable form
\begin{equation}
  \hat V_{\mathrm{nl}}^A
  =\sum_{l m}\sum_{ij}
   |p_{i,lm}^{A}\rangle h_{ij}^{A,l}
   \langle p_{j,lm}^{A}|
  \label{eq:gth-nonlocal-operator}
\end{equation}
이다. AO matrix element를 projector overlap
\begin{equation}
  P_{i,lm;\mu}^{A\vk}
  =\langle p_{i,lm}^{A}|\phi_{\mu\vk}\rangle
  \label{eq:gth-projector-overlap}
\end{equation}
로 쓰면
\begin{equation}
  V_{\mu\nu,\mathrm{nl}}^{\vk}
  =\sum_{A,lm,ij}
   P_{i,lm;\mu}^{A\vk *}h_{ij}^{A,l}
   P_{j,lm;\nu}^{A\vk}.
  \label{eq:gth-nonlocal-matrix}
\end{equation}
Product rule로
\begin{align}
  V_{\mu\nu,\mathrm{nl}}^{\vk,Ax}
  =\sum_{l m i j}
  \bigg[&P_{i,lm;\mu}^{A\vk,Ax *}h_{ij}^{A,l}
         P_{j,lm;\nu}^{A\vk}
  \notag\\
       +&P_{i,lm;\mu}^{A\vk *}h_{ij}^{A,l}
         P_{j,lm;\nu}^{A\vk,Ax}\bigg]
  +\text{AO-center terms}
  \label{eq:gth-nonlocal-derivative}
\end{align}
를 얻는다. Projector radial table $h_{ij}^{A,l}$ 자체는 원자와 함께 rigid하게 이동하므로
핵 좌표에 대한 parameter derivative는 없다.

\subsection{AO-row convention과 full matrix}

Integral engine은 흔히 AO row center에 대한 derivative matrix
$h_{\mu\nu}^{(1)}$를 반환한다. Atom $A$의 full matrix derivative를 만들려면
\begin{equation}
  h^{Ax}=h_{\mathrm{row},A}^{x}
       +h_{\mathrm{row},A}^{x\dagger}
       +h_{\mathrm{operator},A}^{x}
  \label{eq:hcore-full-derivative}
\end{equation}
처럼 bra/ket와 operator-center ownership을 합쳐야 한다. Complex k-point에서 두 번째
항은 단순 transpose가 아니라 conjugate transpose다.

Overlap derivative도 같은 AO-center logic으로
\begin{equation}
  S^{Ax}=S_{\mathrm{row},A}^{x}+S_{\mathrm{row},A}^{x\dagger}
  \label{eq:overlap-full-derivative}
\end{equation}
를 만족한다. $h^{Ax}$와 $S^{Ax}$는 explicit target gradient뿐 아니라 source HF CPHF
source 식~\eqref{eq:cphf-source}에도 사용된다.

\subsection{Pseudopotential 일관성}

Source density를 HF로 만들었다는 이유로 ``HF pseudopotential''이라는 별도의 analytic
형식이 생기지는 않는다. GTH-HF와 같은 label은 pseudopotential parameter set을 뜻한다.
Source HF와 target DFA energy/gradient는 동일한 GTH parameter와 valence ionic charge를
사용해야 한다. 서로 다른 pseudopotential로 source와 target을 계산하면 식
\eqref{eq:scope-dc-energy}의 한 energy surface를 미분하는 문제가 아니게 된다.

검증에서는 local, nonlocal, kinetic/overlap, ion--ion component를 분리해 FD와 비교한 뒤
전체 hcore derivative를 확인한다. 특히 nonlocal projector의 angular channel이 늘어나는
basis/pseudopotential에서 atom slice와 conjugation 오류가 잘 드러난다.
